% Mathematical companion to doc/designs/ipopt_outer_bound_design.md
% Build:  pdflatex ipopt_outer_bound_certificate.tex
\documentclass[11pt]{article}

\usepackage[margin=1.1in]{geometry}
\usepackage{amsmath,amssymb,amsthm}
\usepackage{booktabs}
\usepackage[hidelinks]{hyperref}

\theoremstyle{plain}
\newtheorem{lemma}{Lemma}
\newtheorem{theorem}[lemma]{Theorem}
\newtheorem{corollary}[lemma]{Corollary}
\newtheorem{proposition}[lemma]{Proposition}
\theoremstyle{remark}
\newtheorem{remark}[lemma]{Remark}

\newcommand{\R}{\mathbb{R}}
\newcommand{\W}{W}
\newcommand{\vhat}{\hat{v}}
\newcommand{\qhat}{\hat{q}}
\newcommand{\OPT}{\mathrm{OPT}}

\title{A certified outer bound for convex NLP scenario subproblems}
\author{mpi-sppy \\ \small companion to \texttt{doc/designs/ipopt\_outer\_bound\_design.md}}
\date{}

\begin{document}
\maketitle

\begin{abstract}
The usual device for getting outer bounds in mip-sppy is to read the dual bound off the subproblem solver, which fails for Ipopt, which reports no dual bound.  This note
derives the bound that the \texttt{ipopt\_outer\_bound} spoke computes instead.
The construction is Lagrangian weak duality combined with a tangent-plane
underestimator minimised in closed form over the variable box.  \emph{None of the
mathematics is new}: Theorem~\ref{thm:main} is the Frank--Wolfe duality gap on a
box, and Section~\ref{sec:literature} places each result against the
literature.  The derivation is written out because the hypotheses are what
matter in practice, and one of them is easy to get backwards.  Its defining
property is that it requires \emph{no} assumption about the accuracy of the
solve: a truncated or sloppy solve yields a loose bound rather than an invalid
one.  Convexity, by contrast, is genuinely load-bearing, and
Section~\ref{sec:canonical} isolates the one place where it is easy to get
backwards.
\end{abstract}

\section{Setting}

We assume minimization. 
Fix a scenario $s$ with probability $p_s$, $\sum_s p_s = 1$.  Let
$v \in \R^{n}$ collect \emph{all} variables of the scenario subproblem --- the
nonanticipative variables $x(v)$ together with the recourse variables --- and let
\[
  B \;=\; \prod_{i=1}^{n} [\,l_i, u_i\,]
\]
be the box of variable bounds.  With hub weights $\W_s$, the subproblem an outer bound
spoke solves is
\begin{equation}\label{eq:sub}
  L_s(\W_s) \;=\; \min_{v} \;\bigl\{\, f_s(v) + \W_s^{\top} x(v)
     \;:\; g_s(v) \le 0,\; h_s(v) = 0,\; v \in B \,\bigr\},
\end{equation}
with $f_s : \R^n \to \R$, $g_s : \R^n \to \R^{m}$ and $h_s : \R^n \to \R^{k}$.
This is exactly the problem the ordinary Lagrangian spoke builds, with the
proximal term switched off.

\paragraph{Why the returned objective value is useless here.}
For a minimisation, the value the solver reports is the objective evaluated at
the point it stopped at.  Any such value is $\ge$ the minimum, so it is an
\emph{inner} bound: the wrong direction.  Nor does Ipopt's convergence tolerance
close the gap, since it is a relative KKT-residual tolerance rather than a bound
on the objective error, and the two error sources point opposite ways ---
optimality error stops slightly above the optimum, while constraint violation
can place the iterate slightly below it.

\section{Weak duality}

\begin{lemma}[Weak duality]\label{lem:weak}
Let $\lambda \in \R^{m}_{+}$ and $\mu \in \R^{k}$ be arbitrary.  Define
\begin{align}
  \varphi_s(v) &\;=\; f_s(v) + \W_s^{\top} x(v)
                     + \lambda^{\top} g_s(v) + \mu^{\top} h_s(v),
     \label{eq:phi}\\
  q_s(\lambda,\mu) &\;=\; \inf_{v \in B} \varphi_s(v).
     \label{eq:q}
\end{align}
Then $q_s(\lambda,\mu) \le L_s(\W_s)$.
\end{lemma}

\begin{proof}
Let $v$ be any point feasible for \eqref{eq:sub}.  Then $v \in B$, so
$q_s(\lambda,\mu) \le \varphi_s(v)$.  Moreover $g_s(v) \le 0$ and $\lambda \ge 0$
give $\lambda^{\top} g_s(v) \le 0$, and $h_s(v) = 0$ gives
$\mu^{\top} h_s(v) = 0$; hence
$\varphi_s(v) \le f_s(v) + \W_s^{\top} x(v)$.  Chaining the two inequalities,
$q_s(\lambda,\mu) \le f_s(v) + \W_s^{\top} x(v)$ for every feasible $v$.  Taking
the infimum over feasible $v$ gives the claim.
\end{proof}

Note what Lemma~\ref{lem:weak} does \emph{not} require: $\lambda$ and $\mu$ need
not be optimal, or even good.  Any $\lambda \ge 0$ and any $\mu$ will do.  This
is the property that makes the whole approach robust, and it is why a mistaken
sign convention or a stale multiplier can only cost tightness
(Remark~\ref{rem:robust}).

\paragraph{The trap.}
Lemma~\ref{lem:weak} is not yet usable, because $q_s$ is defined by an
\emph{infimum}.  Handing $\varphi_s$ to a nonlinear solver returns a
\emph{point}, and the value there is $\ge$ the infimum --- an upper bound on
$q_s$, which is the wrong direction again.  A second NLP solve therefore does
not by itself certify anything.

\section{The box underestimator}

\begin{theorem}[Certified bound]\label{thm:main}
Suppose $\varphi_s$ is convex on an open set containing $B$ and differentiable
at $\vhat$.  Define
\begin{equation}\label{eq:qhat}
  \qhat_s(\vhat) \;=\; \varphi_s(\vhat)
    \;+\; \sum_{i=1}^{n} \; \min_{t \in [l_i,u_i]}
          \partial_i \varphi_s(\vhat)\,\bigl(t - \vhat_i\bigr).
\end{equation}
Then $\qhat_s(\vhat) \le q_s(\lambda,\mu) \le L_s(\W_s)$, and each term of the
sum is available in closed form:
\begin{equation}\label{eq:closed}
  \min_{t \in [l_i,u_i]} \partial_i \varphi_s(\vhat)\,(t - \vhat_i)
  \;=\;
  \begin{cases}
    \partial_i \varphi_s(\vhat)\,(l_i - \vhat_i), & \partial_i \varphi_s(\vhat) > 0,\\[2pt]
    \partial_i \varphi_s(\vhat)\,(u_i - \vhat_i), & \partial_i \varphi_s(\vhat) < 0,\\[2pt]
    0, & \partial_i \varphi_s(\vhat) = 0.
  \end{cases}
\end{equation}
\end{theorem}

\begin{proof}
Convexity and differentiability at $\vhat$ give the gradient inequality
\[
  \varphi_s(v) \;\ge\; \varphi_s(\vhat)
      + \nabla \varphi_s(\vhat)^{\top} (v - \vhat)
  \qquad \text{for all } v .
\]
Minimising both sides over $v \in B$ preserves the inequality, and the right-hand
side is separable across coordinates because $B$ is a box, which yields
\eqref{eq:qhat}.  Each one-dimensional problem minimises a linear function over
an interval, so the minimum is attained at an endpoint, giving
\eqref{eq:closed}.  The final inequality is Lemma~\ref{lem:weak}.
\end{proof}

\begin{corollary}\label{cor:anypoint}
$\qhat_s(\vhat) \le L_s(\W_s)$ for \emph{any} $\vhat$ at which $\varphi_s$ is
differentiable, \emph{any} $\lambda \ge 0$ and \emph{any} $\mu$.  In particular
$\vhat$ need not be optimal, and need not even be feasible for \eqref{eq:sub}.
\end{corollary}

Corollary~\ref{cor:anypoint} is the whole point.  The certificate consumes the
iterate the solver happened to stop at and the multipliers it happened to
report, and returns a valid bound regardless.  Computationally it costs one
gradient evaluation and a loop over the variables --- no second solve.

\begin{remark}[Robustness]\label{rem:robust}
Because Corollary~\ref{cor:anypoint} holds for arbitrary $\lambda \ge 0$ and
$\mu$, an error in the sign convention used to recover the multipliers, a stale
multiplier, or a multiplier clipped to zero can only make $\qhat_s$ smaller.
Such an error makes the bound loose, never invalid.  Convexity enjoys no such
protection; see Section~\ref{sec:canonical}.
\end{remark}

\section{Exactness at a KKT point}

The correction term in \eqref{eq:qhat} is not merely small in practice: it
vanishes identically at an exact KKT point, so the certificate degrades
gracefully and costs nothing when the solve is good.

\begin{proposition}[The correction vanishes]\label{prop:kkt}
Let $\vhat$ satisfy the KKT conditions of \eqref{eq:sub} with multipliers
$\lambda \ge 0$ for $g_s \le 0$, $\mu$ for $h_s = 0$, and $z_L, z_U \ge 0$ for
the bounds $l - v \le 0$ and $v - u \le 0$.  Then every term of the sum in
\eqref{eq:qhat} is zero, and if in addition $\varphi_s$ is convex then
$\qhat_s(\vhat) = L_s(\W_s)$.
\end{proposition}

\begin{proof}
Stationarity of the full Lagrangian, which is $\varphi_s$ augmented with the
bound terms, reads
\begin{equation}\label{eq:stat}
  \nabla \varphi_s(\vhat) \;=\; z_L - z_U .
\end{equation}
Fix a coordinate $i$ and consider the three possible positions of $\vhat_i$.

If $l_i < \vhat_i < u_i$, complementarity gives $z_{L,i} = z_{U,i} = 0$, so
$\partial_i \varphi_s(\vhat) = 0$ by \eqref{eq:stat} and the term is $0$ by
\eqref{eq:closed}.

If $\vhat_i = l_i$, complementarity gives $z_{U,i} = 0$, so
$\partial_i \varphi_s(\vhat) = z_{L,i} \ge 0$.  When the derivative is strictly
positive \eqref{eq:closed} selects $t = l_i = \vhat_i$, and when it is zero the
term is zero outright; either way the term is $0$.

If $\vhat_i = u_i$, symmetrically $z_{L,i} = 0$ and
$\partial_i \varphi_s(\vhat) = -z_{U,i} \le 0$, so \eqref{eq:closed} selects
$t = u_i = \vhat_i$ and the term is $0$.

Hence $\qhat_s(\vhat) = \varphi_s(\vhat)$.  Complementarity
$\lambda^{\top} g_s(\vhat) = 0$ and feasibility $h_s(\vhat) = 0$ reduce
\eqref{eq:phi} to $\varphi_s(\vhat) = f_s(\vhat) + \W_s^{\top} x(\vhat)$.  Under
convexity the KKT conditions are sufficient for global optimality of
\eqref{eq:sub}, so that value is $L_s(\W_s)$.
\end{proof}

Proposition~\ref{prop:kkt} has a practical corollary worth stating plainly: the
correction term \emph{measures its own inexactness}.  It is zero when the solve
is exact and grows as the solve degrades, so the reported bound loosens smoothly
rather than becoming wrong.

\section{Canonical form, and where convexity is easy to get backwards}
\label{sec:canonical}

Theorem~\ref{thm:main} requires $\varphi_s$ convex, and by \eqref{eq:phi} that
requires each component of $g_s$ to be convex and each component of $h_s$ to be
affine, since $\lambda \ge 0$ but $\mu$ is free in sign.

The subtlety is that $g_s$ is the constraint in \emph{canonical} form
$g_s(v) \le 0$, and putting a $\ge$ row into canonical form negates its body.
Writing $b(v)$ for the constraint body as the user stated it:

\begin{center}
\begin{tabular}{lll}
\toprule
as written & canonical form & requirement on $b$ \\
\midrule
$b(v) \le \beta$              & $g = b - \beta$              & $b$ convex \\
$b(v) \ge \alpha$             & $g = \alpha - b$             & $b$ \textbf{concave} \\
$\alpha \le b(v) \le \beta$   & both rows                    & $b$ affine \\
$b(v) = \beta$                & $h = b - \beta$              & $b$ affine \\
\bottomrule
\end{tabular}
\end{center}

So $x^2 \le 4$ is admissible and $x^2 \ge 1$ is not, though both are written with
a convex body --- and indeed the feasible set of the second is not convex.  This
is the one place in the construction where an entirely ordinary-looking model
silently leaves the hypotheses, and the consequence is not a loose bound but an
invalid one.  A worked instance: for
\[
  \min\; x \quad \text{s.t.} \quad x^2 \ge 1, \quad x \in [0, 1.5],
\]
whose optimum is $1$, taking $\vhat = 0.5$ and $\lambda = 1$ in \eqref{eq:qhat}
yields $\qhat = 1.25 > 1$.

Of the four rows above, three are decidable --- affineness is a question about
polynomial degree --- and the implementation enforces them.  Convexity or
concavity of a one-sided nonlinear body is not decidable, and remains a user
assertion.

\section{Aggregating across scenarios}

The per-scenario bounds are combined by the usual probability-weighted sum, and
the condition under which that sum is itself a bound deserves care.

\begin{proposition}[Aggregation]\label{prop:agg}
Let $\OPT$ be the optimal value of the original stochastic program.  If the
weights satisfy $\sum_s p_s \W_s = 0$, then
$\sum_s p_s\, L_s(\W_s) \le \OPT$.
\end{proposition}

\begin{proof}
Let $x^{*}$ together with the scenario recourse decisions be optimal for the
original problem, so that $\OPT = \sum_s p_s f_s(v_s^{*})$ where $v_s^{*}$ is the
optimal decision in scenario $s$ and $x(v_s^{*}) = x^{*}$ for every $s$ by
nonanticipativity.  Each $v_s^{*}$ is feasible for subproblem \eqref{eq:sub}, so
$L_s(\W_s) \le f_s(v_s^{*}) + \W_s^{\top} x^{*}$.  Multiplying by $p_s$ and
summing,
\[
  \sum_s p_s L_s(\W_s)
  \;\le\; \sum_s p_s f_s(v_s^{*}) + \Bigl( \sum_s p_s \W_s \Bigr)^{\!\top} x^{*}
  \;=\; \OPT + 0 . \qedhere
\]
\end{proof}

\begin{remark}[Weights must share a generation]\label{rem:mixed}
The hypothesis $\sum_s p_s \W_s = 0$ is a \emph{joint} condition on the whole
weight vector, and progressive hedging maintains it at each iteration.  It is
not preserved by mixing: if $\W'_s$ takes scenario $s$'s weights from one
iteration and scenario $t$'s from another, then in general
$\sum_s p_s \W'_s \ne 0$, the proof above leaves the residual term
$\bigl(\sum_s p_s \W'_s\bigr)^{\top} x^{*}$ in place, and the resulting number
is not a bound at all.  Nothing downstream can detect this, which is why a
scenario whose certificate is unavailable must make the whole expectation
unavailable rather than contribute a stale value.
\end{remark}

\section{Practical consequences}

\paragraph{Tightness is governed by the box.}
By \eqref{eq:closed} the shortfall of $\qhat_s$ below $\varphi_s(\vhat)$ is
\[
  \sum_{i=1}^{n} \bigl| \partial_i \varphi_s(\vhat) \bigr| \cdot
  \bigl( \text{distance from } \vhat_i \text{ to the far end of } [l_i,u_i] \bigr),
\]
so the bound is only as tight as the variable bounds are.  Gradient components
at a converged solve are small but generically nonzero, so a box of width $10$
costs little while a box of width $10^{10}$ can cost a substantial fraction of
the objective.  Tightening bounds before certifying --- for instance by
feasibility-based bounds tightening, which shrinks $B$ without removing feasible
points and therefore only raises $\qhat_s$ --- is the single most effective
lever on quality.

\paragraph{Unbounded variables.}
If $l_i = -\infty$ and $\partial_i \varphi_s(\vhat) > 0$, or $u_i = +\infty$ and
$\partial_i \varphi_s(\vhat) < 0$, the corresponding term in \eqref{eq:closed} is
$-\infty$ and no finite bound is available.  Whether this occurs depends on the
model and on how converged the solve is: at an exact KKT point the offending
component is zero by Proposition~\ref{prop:kkt}, but away from one it generally
is not.  The implementation reports ``no bound'' in that case rather than
$-\infty$, so that Remark~\ref{rem:mixed} is respected.

\paragraph{Floating point.}
Nothing in Theorem~\ref{thm:main} requires a tolerance.  A small relative
cushion, reporting $\qhat_s - \varepsilon(1 + |\qhat_s|)$, is nevertheless cheap
insurance against last-bit effects in the gradient and the sum; it is hygiene,
not part of the argument.

\section{Relation to known results}\label{sec:literature}

Nothing in Sections~2--6 is new, and it is worth saying which wheel each result
re-uses.

\paragraph{Lemma~\ref{lem:weak}} is textbook Lagrangian weak duality, stated for
an arbitrary multiplier pair rather than an optimal one.

\paragraph{Theorem~\ref{thm:main}} is the \emph{Frank--Wolfe duality gap},
sometimes called the Wolfe gap or the linearisation gap.  For a convex $f$ on a
compact convex $D$, evaluating the linear minimisation oracle at $x$ gives
\[
  f(x) + \min_{y \in D} \nabla f(x)^{\top}(y - x) \;\le\; \min_{D} f ,
\]
which is \eqref{eq:qhat} with $D = B$.  The bound is as old as the Frank--Wolfe
algorithm itself \cite{frankwolfe1956}; its modern reading as a \emph{certificate}
computable at any iterate, rather than as a step-size device, is the framing in
\cite{jaggi2013}.  The only specialisation here is that a box makes the oracle
separable and closed-form, so the certificate costs one gradient evaluation and
no optimisation at all.

\paragraph{Proposition~\ref{prop:agg}} is the standard argument by which
progressive hedging weights yield a lower bound on the optimal value, as in
\cite{gade2016}.

\paragraph{The combination} of progressive hedging with Frank--Wolfe machinery to
obtain Lagrangian dual bounds is itself established \cite{boland2018}, and
mpi-sppy already ships an FWPH cylinder on that basis.

What is specific here is therefore an assembly rather than a result: the
multipliers are the ones Ipopt already returns, the linearisation point is the
iterate it already stopped at, and the box is the model's own variable bounds ---
so a valid outer bound is obtained from the solve the spoke was going to perform
anyway, with no second optimisation and no assumption that the first one
converged.  The wider question of extracting valid dual bounds from an inexact
oracle is treated at length in the inexact bundle-method literature, which this
note does not attempt to survey.

\section{Correspondence with the implementation}

\begin{center}
\begin{tabular}{ll}
\toprule
object here & in the code \\
\midrule
$\varphi_s$ in \eqref{eq:phi}      & \texttt{\_lagrangian\_expression} \\
$\qhat_s$ in \eqref{eq:qhat}       & \texttt{certified\_lower\_bound} \\
decidable hypotheses               & \texttt{check\_model\_is\_certifiable} \\
box tightening                     & \texttt{unbounded\_variables} \\
Proposition~\ref{prop:agg} sum     & \texttt{SPOpt.Ebound} \\
Remark~\ref{rem:mixed} safeguard   & a \texttt{None} bound, propagated collectively \\
\bottomrule
\end{tabular}
\end{center}

\noindent
The first four live in \texttt{mpisppy/utils/dual\_certificate.py}; the spoke
that drives them is \texttt{mpisppy/cylinders/ipopt\_outer\_bound.py}.

\begin{thebibliography}{9}

\bibitem{frankwolfe1956}
Frank, M. and Wolfe, P.:
\newblock An algorithm for quadratic programming.
\newblock \emph{Naval Research Logistics Quarterly} \textbf{3}(1--2), 95--110 (1956).

\bibitem{jaggi2013}
Jaggi, M.:
\newblock Revisiting Frank--Wolfe: projection-free sparse convex optimization.
\newblock \emph{Proceedings of the 30th International Conference on Machine
Learning (ICML)}, 427--435 (2013).

\bibitem{gade2016}
Gade, D., Hackebeil, G., Ryan, S.M., Watson, J.-P., Wets, R.J.-B. and
Woodruff, D.L.:
\newblock Obtaining lower bounds from the progressive hedging algorithm for
stochastic mixed-integer programs.
\newblock \emph{Mathematical Programming} \textbf{157}(1), 47--67 (2016).

\bibitem{boland2018}
Boland, N., Christiansen, J., Dandurand, B., Eberhard, A., Linderoth, J.,
Luedtke, J. and Oliveira, F.:
\newblock Combining progressive hedging with a Frank--Wolfe method to compute
Lagrangian dual bounds in stochastic mixed-integer programming.
\newblock \emph{SIAM Journal on Optimization} \textbf{28}(2), 1312--1336 (2018).

\end{thebibliography}

\end{document}
